\section{Two-Layer Validation and Recovery}
\label{sec:framework}
\begin{figure}[t]
\centering
\begin{tikzpicture}[node distance=5mm,font=\scriptsize,
box/.style={draw,rounded corners,align=center,text width=6.5cm,minimum height=7mm},
arr/.style={-{Stealth[length=2mm]}}]
\node[box,fill=blue!7] (train) {Enrolled clients train on private training data};
\node[box,below=of train,fill=blue!7] (vote) {Layer 1: eligible committee evaluates candidates\\on separate private validation data; quorum approves updates};
\node[box,below=of vote,fill=green!7] (commit) {Off-chain aggregation; on-chain model commitment\\Account-authenticated votes and stake accounting};
\node[box,below=of commit,fill=orange!9] (monitor) {Layer 2: fixed-panel windowed quality monitoring};
\node[box,below=of monitor,fill=orange!9] (recover) {Flag: append recovery and penalty events\\Restore historical off-chain checkpoint; retain all transactions};
\draw[arr] (train)--(vote);\draw[arr] (vote)--(commit);
\draw[arr] (commit)--(monitor);\draw[arr] (monitor)--(recover);
\end{tikzpicture}
\caption{System workflow. Raw examples remain local. The contracts enforce
protocol transitions; they do not verify private quality measurements.}
\Description{Five vertically connected steps: client training, private committee validation, model commitment, quality monitoring, and append-only recovery.}
\label{fig:system}
\end{figure}

\subsection{Layer 1: Eligibility and Validation}
For candidate client $k$, let $E$ be the number of eligible validators after
excluding $k$. Enrollment, minimum stake, and the configured trust threshold
determine eligibility. The committee size and approval threshold are
\begin{equation}
 C=\min\{E,\max\{3,\lceil0.3E\rceil\}\},\qquad
 q=\lceil2C/3\rceil.
\end{equation}
When $E<3$, validation is unsupported; it is never replaced with a one-member
committee. For example, four enrolled eligible clients give a three-member
committee after excluding the proposer and require two affirmative votes.
Each candidate receives an independently sampled committee with a reproducible
random stream. Committees are sampled independently of participant selection.

Validator $j$ evaluates the candidate and current model on the same private
validation partition. With utility $u_j$, its normalized change is
\begin{equation}
 d_{jk}=\frac{u_j(w+\Delta_k)-u_j(w)}{\max\{|u_j(w)|,10^{-12}\}}.
\end{equation}
The default approval rule is $d_{jk}\geq-0.05$. A candidate requires $q$
distinct affirmative signers; repeated votes cannot increase support.
The median normalized change is recorded diagnostically and is not substituted
for quorum. Accepted updates are aggregated as
\begin{equation}
 w'=w+\frac{\sum_{k\in A} n_k\tau_k\Delta_k}{\sum_{k\in A}n_k\tau_k},
\end{equation}
where $A$ is the approved set, $n_k$ is its training sample count, and $\tau_k$
is recorded trust. Empty or unsupported approval sets produce explicit no-op
rounds. Trust changes follow recorded approval or rejection outcomes, with
additional penalties after recovery. They do not use ground-truth attacker labels.
Trust starts at $0.5$; validator eligibility requires $\tau_k\geq0.45$, while
candidate training stops below the floor $0.15$. Approval adds $0.02$ trust;
rejection or a recovery-attributed penalty subtracts $0.15$, clipped to $[0,1]$.
Initial stake is 20 credits, eligibility requires 10, and an attributed penalty
deducts five. These are prespecified starting values rather than theoretically
optimal thresholds. Eligibility sensitivity includes $0.30$, $0.45$, and $0.50$;
its outcomes are needed to justify a final deployment choice.

\subsection{Layer 2: Monitoring a Fixed Panel}
A panel is fixed for a window of five observations, including the baseline
before its first candidate round. Panel members evaluate consecutive versions
on their private validation partitions. Keeping panel membership constant within
the window reduces score changes caused solely by resampling validators.
Let $s_i$ be the median panel utility at observation $i$. With $W=5$, the monitor
compares the latest score with the mean of preceding window scores and fits
a least-squares linear slope:
\begin{equation}
 D=\frac{\overline{s}_{1:W-1}-s_W}{\max\{|\overline{s}_{1:W-1}|,10^{-12}\}}.
\end{equation}
A relative drop above $0.05$ or a slope below $-0.005$ triggers recovery.
An adjacent-round detector is evaluated separately. Thresholds are prespecified
and their sensitivity is measured using development seeds. Non-IID data and benign
drift can trigger both monitors, so detection does not prove an attack occurred.

\subsection{Append-Only Recovery and Accountability}
Every model version has a commitment, parent pointer, and off-chain checkpoint.
After a flagged window, recovery restores the checkpoint preceding that window.
The contract appends the invalidated interval, restored version, monitoring-panel
attestations, and penalty events. Earlier commitments remain retrievable.
The restored checkpoint is historical and may contain older poisoning.

Only affirmative signers of commitments in the invalidated interval are eligible
for penalties. A signer is penalized at most once per recovery event. Rejected
votes and nonparticipants are not penalized. Stake and trust changes persist after
model restoration; the detector history resets and a new panel is selected.
Because honest signers can approve harmful updates or a monitor can misfire,
wrongful penalties are a primary outcome rather than an omitted corner case.
